\documentclass[12pt,a4paper]{article}

% ==========================================
% 1. ENCODAGE ET POLICES (FONTS)
% ==========================================
\usepackage[T1]{fontenc}   % Encodage des polices de sortie pour la gestion des accents
\usepackage{times}        % Police Times New Roman pour le texte
\usepackage{lmodern}      % Police Latin Modern améliorée
\usepackage{pifont}       % Symboles et caractères spéciaux (ex: \ding{45})

% ==========================================
% 2. DISPOSITION DE LA PAGE ET EN-TÊTES/PIEDS DE PAGE
% ==========================================
\usepackage[left=1cm, right=1cm, top=1cm, bottom=1.7cm]{geometry} % Configuration des marges
\usepackage{fancyhdr}     % Personnalisation des en-têtes et pieds de page
\usepackage{lastpage}     % Accès au numéro de la dernière page pour "page X/Y"
\usepackage{multicol}     % Gestion de la mise en page sur plusieurs colonnes

% Configuration du style des pieds de page
\pagestyle{fancy}
\fancyhf{} % Réinitialise les en-têtes et pieds de page
\renewcommand{\headrulewidth}{0pt}  % Pas de ligne sous l'en-tête
\renewcommand{\footrulewidth}{1pt}  % Ligne au-dessus du pied de page
\renewcommand{\footruleskip}{10pt}  % Espacement de la ligne de pied de page

%\fancyfoot[L]{\color{blue}\ding{45}\ \textbf{Prof: M. Diatta}} % Pied de page gauche
%\fancyfoot[R]{\color{blue}\ding{45}\ \textbf{2026}}        % Pied de page droit
%\cfoot{\bf\thepage\ / \pageref{LastPage}}                 % Numérotation centrale (X/Y)

% ==========================================
% 3. MATHÉMATIQUES ET SYSTÈMES
% ==========================================
\usepackage{amsmath}     % Environnements et commandes mathématiques de base
\usepackage{amssymb}     % Symboles mathématiques avancés
\usepackage{mathrsfs}    % Police calligraphique mathématique (\mathscr)

% Force l'affichage grand format (displaystyle) pour les formules intégrées au texte
\everymath{\displaystyle}

% ==========================================
% 4. TABLEAUX ET FIGURES
% ==========================================
\usepackage{array}       % Amélioration du traitement des tableaux
\usepackage{multirow}    % Fusion de cellules sur plusieurs lignes
\usepackage{float}        % Contrôle avancé du placement des figures ([H])
\usepackage{varwidth}    % Ajustement automatique de la largeur des boîtes

% ==========================================
% 5. LISTES ET PUCES PERSONNALISÉES
% ==========================================
\usepackage{enumitem}    % Personnalisation avancée des listes à puces et numérotées

% Styles TikZ pour puces circulaires et carrées
\newcommand\circitem[1]{%
  \tikz[baseline=(char.base)]{
    \node[circle, draw=gray, fill=red!55, minimum size=1.2em, inner sep=0] (char) {#1};
  }%
}
\newcommand\boxitem[1]{%
  \tikz[baseline=(char.base)]{
    \node[fill=cyan, minimum size=1.2em, inner sep=0] (char) {#1};
  }%
}

% Application aux niveaux 1 et 2 des listes numérotées
\setlist[enumerate,1]{label=\protect\circitem{\arabic*}}
\setlist[enumerate,2]{label=\protect\boxitem{\alph*}}

% ==========================================
% 6. GRAPHISMES, TIKZ ET TABLEAUX DE VARIATION
% ==========================================
\usepackage{tikz}        % Moteur de dessin vectoriel
\usetikzlibrary{patterns, trees} % Extensions TikZ :
                                 % - patterns : motifs de hachures et remplissages
                                 % - trees : réalisation d'arbres de dénombrement/probabilités

\usepackage{tkz-tab}     % Tableaux de variations et de signes mathématiques

% ==========================================
% 7. BOÎTES ET ENCADRÉS (TCOLORBOX)
% ==========================================
\usepackage[most]{tcolorbox} % Boîtes encadrées stylisées
\tcbuselibrary{skins, hooks} % Extensions graphiques (skins) et accroches d'actions (hooks)

% Environnement de boîte personnalisée pour les exemples
\newtcolorbox{exa}[2][]{
  enhanced, breakable,
  before skip=2mm, after skip=5mm,
  colback=yellow!20!white, colframe=black!20!blue, boxrule=0.5mm,
  attach boxed title to top left={xshift=0.6cm, yshift*=1mm-\tcboxedtitleheight},
  fonttitle=\bfseries,
  title={#2}, #1,
  boxed title style={
    frame code={
      \path[fill=tcbcolback!30!black]
        ([yshift=-1mm, xshift=-1mm]frame.north west)
        arc[start angle=0, end angle=180, radius=1mm]
        ([yshift=-1mm, xshift=1mm]frame.north east)
        arc[start angle=180, end angle=0, radius=1mm];
      \path[left color=tcbcolback!60!black, right color=tcbcolback!60!black, middle color=tcbcolback!80!black]
        ([xshift=-2mm]frame.north west) -- ([xshift=2mm]frame.north east)
        [rounded corners=1mm]-- ([xshift=1mm, yshift=-1mm]frame.north east)
        -- (frame.south east) -- (frame.south west)
        -- ([xshift=-1mm, yshift=-1mm]frame.north west)
        [sharp corners]-- cycle;
    },
    interior engine=empty,
  },
  interior style={top color=yellow!5}
}

% ==========================================
% 8. ARRIÈRE-PLAN ET FILIGRANE (WATERMARK)
% ==========================================
\usepackage{eso-pic}     % Ajout d'éléments en arrière-plan des pages

% Filigrane "PGB" centré sur chaque page
\AddToShipoutPicture{
  \AtTextCenter{%
    \makebox[0pt]{\rotatebox{50}{\textcolor[gray]{0.6}{\fontsize{1.5cm}{1.5cm}\selectfont maths-mes-maths}}}
  }
}

% ==========================================
% 9. COMPTEURS ET COMMANDES DE COURS
% ==========================================
% Compteur et commande pour la numérotation des exercices
\newcounter{exercice}
\renewcommand{\theexercice}{\arabic{exercice}}
\newcommand{\exo}{\refstepcounter{exercice}\textbf{Exercice \theexercice}\ }

% ==========================================
% 10. HYPERLIENS ET INTERACTIVITÉ
% ==========================================
% Chargé en tout dernier pour éviter les conflits de redéfinition de commandes
\usepackage[colorlinks=true, linkcolor=blue, urlcolor=blue, citecolor=blue]{hyperref}
\begin{document}
\renewcommand{\arraystretch}{1.5}
\renewcommand{\arrayrulewidth}{1.2pt}
\begin{tcolorbox}[
    colback=white,
    colframe=blue!60!black,
    boxrule=0.8pt,
    arc=3pt,
    left=3mm, right=3mm, top=2.5mm, bottom=2.5mm
]
    \begin{center}
        \textbf{\large Lycée de Diattacounda} \hfill \textbf{A.S. : 2026/2027} \\[1mm]
        \textbf{Mathématiques} $\bullet$ \textbf{Niveau : TS2} \hfill \textbf{Date : 13/11/2026} \\[2mm]
        \rule{0.8\linewidth}{0.4pt} \\[1.5mm]
        {\Large\bfseries\sffamily TD : Nombres Complexes}
    \end{center}
\end{tcolorbox}
\vspace{4mm}

\fbox{\textbf{\exo}}

Dans chacun des cas suivants, déterminer l’écriture complexe de la similitude directe de centre \( \Omega \), de rapport \( k \) et d’angle \( \theta \).

\begin{multicols}{2}
\setlength{\columnseprule}{0.1mm} % La largeur de la ligne verticale entre les colonnes
\begin{enumerate}
    \item \( \Omega = O \), \( k = 2 \) et \( \theta = \dfrac{\pi}{6} \).
    
    \item \( \Omega \begin{pmatrix} 1 \\ 0 \end{pmatrix} \), \( k = \sqrt{2} \) et \( \theta = \dfrac{\pi}{2} \).
    
    \item \( \Omega \begin{pmatrix} 2 \\ -1 \end{pmatrix} \), \( k = 1 \) et \( \theta = \dfrac{5\pi}{6} \).
    
    \item \( \Omega (-2 + i \frac{1}{2}) \), \( k = \dfrac{1}{2} \), \( \theta = \dfrac{-\pi}{3} \).
\end{enumerate}
\end{multicols}

\fbox{\textbf{\exo}}

Dans chacun des cas suivants, déterminer la nature et les éléments caractéristiques de l’application \( f \) du plan dans lui-même, qui au point \( M \) d’affixe \( z \) associe le point \( M' \) d’affixe \( z' \) :

\begin{multicols}{3}
\setlength{\columnseprule}{0.1mm} % La largeur de la ligne verticale entre les colonnes
\begin{enumerate}
    \item \( z' = (\sqrt{3} + i)z \)
    \item \( z' = -2z + i \)
    \item \( z' = -z + 2i \)
    \item \( z' = \dfrac{3 + i\sqrt{3}}{4}z + \dfrac{1 - i\sqrt{3}}{2} \)
    \item \( z' = iz + 1 \)
    \item \( z' = z + 3 - i \)
\end{enumerate}
\end{multicols}

\fbox{\textbf{\exo}}

Dans chacun des cas suivants, déterminer l’écriture complexe de la similitude directe \( s \) définie par :  
\( S(A) = A' \) et \( S(B) = B' \).

\begin{multicols}{2}
\setlength{\columnseprule}{0.1mm} % La largeur de la ligne verticale entre les colonnes
\begin{enumerate}
    \item \( A(3+2i), A'(3), B(1) \) et \( B'(i) \);
    \item \( A(2+i), A'(3+2i), B(2) \) et \( B'(3i) \).
\end{enumerate}
\end{multicols}

\fbox{\textbf{\exo}}

Soit \( s \) la similitude directe d’écriture complexe :  \( z' = 3iz - 9 - 3i. \)

\begin{enumerate}
    \item Déterminer les éléments caractéristiques de \( s \).
    \item Déterminer l’expression analytique de \( s \).
    \item Déterminer l’image par \( s \) :
    \begin{enumerate}
        \item du cercle de centre \( K(1 - 3i) \) et de rayon 1;
        \item de la droite \( (\mathscr{D}) \) d’équation : \( x = 1 \);
        \item de la droite passant par le point \( A(5;3) \) et dont un vecteur directeur est \( \vec{u}(-1;3) \).
    \end{enumerate}
\end{enumerate}

\fbox{\textbf{\exo}}

Soit \( A, B \) et \( C \) les points d’affixes respectives :  
\( i, 1 + i \) et \( 2 + 2i \).

\begin{enumerate}
    \item Déterminer l’affixe du barycentre \( G \) des points \( A, B \) et \( C \) affectés respectivement des coefficients \( 2, -2 \) et \( 1 \).
    \item Démontrer que la similitude directe \( s \), qui transforme \( A \) en \( B \) et \( B \) en \( C \), a pour centre le point \( G \).
    \item Déterminer l’angle et le rapport de \( s \).
\end{enumerate}

\fbox{\textbf{\exo}} (BAC)

On considère les applications \( T_1 \) et \( T_2 \) dont les écritures complexes sont :  
\( T_1 : z_1 = (\sqrt{3} + i)z \) et \( T_2 : z_2 = (1 - i\sqrt{3})z + 3 \).

\begin{enumerate}
    \item Déterminer la nature et les éléments caractéristiques de \( T_1 \) et \( T_2 \).
    \item Déterminer l’écriture complexe, la nature et les éléments caractéristiques de \( T_2 \circ T_1 \).
    \item Démontrer qu’il existe un seul point \( K \) tel que \( T_1(K) = T_2(K) \). Soit \( L = T_1(K) \).  
          Calculer les affixes des points \( K \) et \( L \).
    \item Démontrer que le point \( L \) est invariable pour chacune des applications \( T_2 \circ T_1^{-1} \) et \( T_1 \circ T_2^{-1} \).
    \item Déterminer l’écriture complexe de chacune des applications \( T_2 \circ T_1^{-1} \) et \( T_1 \circ T_2^{-1} \).  
          Préciser leurs natures et éléments caractéristiques.
\end{enumerate}

\fbox{\textbf{\exo}}

Soit \( S_1 \) l’application qui à tout point \( M(x,y) \) associe \( M'(x',y') \) définie par :

\[
\begin{cases}
    x' = -\frac{1}{2}x - \frac{\sqrt{3}}{2}y + \frac{3}{2} \\
    y' = \frac{\sqrt{3}}{2}x - \frac{1}{2}y - \frac{\sqrt{3}}{2}
\end{cases}
\]

\begin{enumerate}
    \item Déterminer l’écriture complexe de \( S_1 \).
    \item En déduire la nature et les éléments caractéristiques de \( S_1 \).
    \item Soient les points \( A(1) \) et \( B(-1) \). Déterminer l’écriture complexe de la similitude directe  
          \( S_2 \) telle que \( S_2(A) = O \) et \( S_2(B) = B \) puis caractériser \( S_2 \).
    \item On pose \( S = S_1 \circ S_2 \). Déterminer l’écriture complexe de \( S \).
    \item Déterminer l’image par \( S \) :
    \begin{enumerate}
        \item de la droite \( (\mathscr{D}) \) d’équation \( 2x + y - 1 = 0 \);
        \item du cercle \( (\mathscr{C}) \) d’équation \( (x + 3)^2 + (y - 2)^2 = 4 \).
    \end{enumerate}
\end{enumerate}

\fbox{\textbf{\exo}} Bac 2024 \hfill \textbf{(05 points)}

\begin{enumerate}

\item Dans le plan complexe muni d’un repère orthonormé direct \( (O; \vec{i}, \vec{j}) \), on considère les points \( A, B \) et \( C \) d’affixes respectives :  
   \( z_A = -3i \), \( z_B = -2 \) et \( z_C = 1 + 2i \).

\begin{enumerate}
    \item Déterminer le module et un argument du quotient  
          \[
          \frac{2z_C - z_B}{z_A - z_B}.
          \]
          \hfill \textbf{(0,75 point)}
          
    \item En déduire la nature du triangle \( ABC \).  
          \hfill \textbf{(0,75 point)}
          
    \item Déterminer l’affixe \( z_D \) du point \( D \) tel que le quadrilatère \( BADC \) soit un carré.  
          \hfill \textbf{(0,5 point)}
          
    \item Montrer que les points \( A, B, C \) et \( D \) appartiennent à un même cercle dont on précisera le centre et le rayon.  
          \hfill \textbf{(0,75 point)}
\end{enumerate}

\item On considère les points \( M \) et \( M' \) d’affixes respectives \( z = x + iy \) et \( z' = x' + iy' \) où \( x, y, x', y' \) sont des réels.

Soit \( S \) l’application du plan dans le plan d’expression analytique : \( \begin{cases}
    x' = x - y + 2 \\
    y' = x + y - 1
\end{cases} \)

\begin{enumerate}
    \item Montrer que l’écriture complexe de \( S \) est :  
          \[
          z' = (1 + i)z + 2 - i.
          \]
          \hfill \textbf{(0,5 point)}
          
    \item Déterminer la nature et les éléments caractéristiques de \( S \).  
          \hfill \textbf{(0,5 point)}
          
    \item Déterminer l’image par \( S \) de la droite \( (D) \) d’équation \( x + y + 1 = 0 \).  
          \hfill \textbf{(0,5 point)}
          
    \item Déterminer l’ensemble des points \( M \) dont l’affixe \( z \) vérifie \( |(1 + i)z + 2 - i| = 2 \).  
          \hfill \textbf{(0,5 point)}
\end{enumerate}
\end{enumerate}
\end{document}

